% ======================================================================
% Section 1 — Introduction
% ======================================================================

\begin{abstract}
\noindent
We prove an exact identity that reduces the maximum loneliness of a Lonely
Runner instance to a finite maximum over cyclic groups: for every speed
vector $V$ of $n$ positive integers,
$M(V)=\max_{\{u,w\}}\tau(u,w)$, where $\tau(u,w)$ is the optimal loneliness
on the pair-sum lattice of denominator $N=v_u+v_w$. The analytic core of the
argument is folklore (it appears in a comment on Tao's 2015 blog post); the
contributions claimed here are the sharpening that sum lattices alone always
suffice, the exact equality, and the lossless reformulation of the Lonely
Runner Conjecture ($\mathrm{LRC}$-$_n$) as the finite statement
$(\mathrm{TAU}$-$_n)$. The reformulation turns the conjecture into a family
of covering problems: a pair certifies iff $n-1$ explicit bad sets
$B_w=\{k\in\Z_N : (n{+}1)\,\wn{wk}<N\}$ fail to cover $\Z_N$. We develop
this covering geometry: an exact size law, a class-group reduction at prime
moduli, a lift law, and complete classifications of the unit coverings at
the prime moduli $7,11,13$ (three bad sets) and $7,13,17,19,37$ (four bad
sets)---at every such prime, covering is governed by a short explicit list
of translate classes in a cyclic group. Empirically, these five primes form
the entire unit-capable \emph{rigid zone} at $n=5$ within the verified range
(primes $\le 150$, composites $\le 111$), a striking but unproved finiteness
statement. On the full corpus of $3{,}712{,}576$ primitive speed sets with
$n=5$ and speeds $\le 56$, proved lemmas close $75.6\%$ of instances
($78.2\%$ including a coincidence-counting lemma), per-modulus verification
closes all of them, and the residual obstruction is localized to a single
named class (mixed order signatures at composite moduli). The proved-coverage
rate declines as the corpus grows ($93.4\%\to78.2\%$ over $V=24\to56$ on the
augmented battery, crossing below the continuation threshold at $V=56$);
the trend does not establish that coverage stabilizes, and its asymptotics
are unresolved. We report the decline, the localization, and the exact
verification record as the paper's empirical findings.
The paper's culminating result is stated as a \emph{conditional}
theorem (Theorem~\ref{thm:SC}): under three named open hypotheses---the
multi-fiber survivor inequality Hq (quantified in this revision; it
subsumes the former dilate-sparsity input H2m as its one-step case),
the distinct-family volume law H1g with its uniform margin, and the
resonance budget H2r---the fiber cascade closes the covering
obstruction at the frontier cells $(1K,T{-}3U)$ of every rung $T\ge 8$
at $N=p^2$ within a depth that is $T$-flat in practice and never
exceeds $T+3$ in the worst case. The reduction chain and every other
input (normal form, reflection fiber, chain-class theorem,
renormalization lemma, lossless equivalence) are proved or
machine-validated, and each constant's provenance---derived, measured,
or fitted, with the data it is fit to---is disclosed in the hypothesis
ledger.
That boundary between the proved and the hypothesized has a structural
type, verified statement-by-statement in \S\ref{sec:scale-status}: every
proved result in this program is a closure, invariance, exhaustiveness,
or machine-verification statement, while every remaining obligation is a
uniform growth-in-$k$ margin over a trivial or measured baseline---a
statement type no turn of the program has proved. This type mismatch is
why the paper stops where it stops. The one candidate instrument
shaped to produce that missing type was a transfer-matrix or automaton
model of the strand census. It was executed as a proof attempt under a
pre-committed success criterion and closed by theorem (for automata
built on the cascade's bijective drift)
(Lemma~\ref{lem:monomial}). The cascade's constrained transfer
matrices are partial permutations with a spectral radius of exactly $0$
or $1$ at every scale. The cyclic product is diagonal, with trace equal
to the survivor count itself, and the missing margin is provably
external to the lossless dynamics. The reformulation is bounded by its
own fidelity: the losslessness that makes it a faithful rewrite is the
property that withholds the decay a margin would require.
This is a structural paper: it does not solve or reduce the Lonely
Runner Conjecture, proves no new cases of it, and claims nothing beyond
the verified ranges stated in the text; the two hypotheses of
Theorem~\ref{thm:SC} are open and are not silently established anywhere.
\end{abstract}

\section{Introduction}\label{sec:intro}

\subsection{The conjecture and its status}

Let $n\ge 2$ and let $V=(v_1,\dots,v_n)$ be distinct positive integer
speeds. Runners start together at the origin of the unit circle
$\T=\R/\Z$ at time $0$; runner $p$ occupies $v_p t \bmod 1$ at time $t$.
Write $\wnc{x}$ for the distance of a real number to the nearest integer.
The \emph{loneliness} of the stationary observer (runner $0$) at time $t$
is
\[
g(t)\;=\;\min_{1\le p\le n}\;\wnc{v_p t},
\qquad
M(V)\;=\;\max_{t\in\T}\,g(t).
\]
The Lonely Runner Conjecture ($\mathrm{LRC}$-$_n$), posed by Wills in 1967
and independently by Cusick in the language of view obstruction
\cite{PS-survey,Wills,Cusick1973}, asserts $M(V)\ge 1/(n{+}1)$ for every $V$ with
$\gcd(V)=1$; by translation symmetry this is the statement that every
runner on the track becomes lonely at some time. The conjecture is known
for at most seven runners (Bohman--Bohman--Holzman and
Barajas--Serra; see the survey \cite{PS-survey} for history and precise
references). The frontier has since moved entirely into computer-assisted
territory: eight runners \cite{Rosenfeld8}, nine and ten runners
\cite{Trakulthongchai}, eleven through thirteen runners
\cite{Sungkawichai}, and fourteen and fifteen runners \cite{Allikvere},
the last with public certificates, resting on the linearly-exponential
checking bounds of Malikiosis--Santos--Schymura \cite{MSS}.

Two neighboring literatures study the fine structure of the problem rather
than new rungs. The \emph{spectrum} line computes exact optimal values
$\mathrm{ML}=p/q$ and their denominators
\cite{Kravitz,FanSun,GiriKravitz,Cordella}, including tight-instance
classifications \cite{Zhang}. The \emph{relations} line localizes
counterexamples and tight instances to low-weight integer hyperplanes in
parameter space \cite{BeckEverett}. The closest formulation in spirit is
Rifford's \cite{Rifford}: he studies the \emph{time} it takes a runner to
become lonely, shows that bounding that time by a number of rounds is
equivalent to a covering problem in dimension $n-2$ on the continuous torus,
formulates a covering conjecture there and proves it for $n\le 6$, and
bounds the one-round gap of loneliness by $1/(2m-1)$. What the two
formulations share is the move from runners to coverings; what separates
them is the domain and the mode. Rifford's covering problem is continuous
and its conjecture is a sufficient condition for a time bound, while the
present reduction is finite and exact---an identity, not an
inequality---reducing the threshold question itself to coverings of the
cyclic groups $\Z_{v_u+v_w}$ by multiplicative bad sets. That finiteness is
what makes the arithmetic visible: nothing corresponding to the rigid zone,
the translate-class classifications, or the kernel world appears in, or is
excluded by, the continuous formulation, and none of our results overlap
his. The present paper is a reformulation paper whose subject is the
geometry of the obstruction itself, at small moduli, in residue language.

\subsection{What this paper does}

The starting point is a classical observation: the optimal time is rational,
and its denominator is the sum or difference of two speeds. In the form
relevant here---``the maximum is attained at an intersection of two sine
curves, so the optimal time is rational with denominator the sum or
difference of two integer velocities''---this appears in a comment on Tao's
blog post of 13 May 2015 \cite{Tao-blog}. We treat that technique as known
folklore and cite it as such. What the folklore does not record, and what
this paper builds on, is the following sharpening and its consequences.

\begin{enumerate}[leftmargin=1.6em,itemsep=2pt]
\item \textbf{The sum-only equality theorem}
  (\S\ref{sec:equality}). At a global optimum with $0<M<\tfrac12$ the
  binding runners can be chosen two-sided, which forces the optimal time
  onto the lattice of a \emph{sum} $v_u+v_w$; difference denominators never
  carry the optimum. On that lattice the two binding runners are exact
  reflections of each other at \emph{every} lattice time
  (Lemma~\ref{lem:mirror}), so the $n$-runner minimum there equals an
  $(n{-}1)$-runner minimum. Together these give the exact identity
  $M(V)=\max_{\{u,w\}}\tau(u,w)$ of Theorem~\ref{thm:equality}---not an
  inequality---and the lossless equivalence of
  Corollary~\ref{cor:tau}: $\mathrm{LRC}$-$_n$ holds iff
  $(\mathrm{TAU}$-$_n)$ does, a \emph{finite} statement about explicit
  cyclic groups $\Z_{v_u+v_w}$. The finitization is exact: no information
  is lost, and nothing about the classical difficulty of the conjecture is
  claimed to become easier.
\item \textbf{The covering framework} (\S\ref{sec:covering}). Pair
  certification becomes a covering question: a pair fails iff its $n-1$ bad
  sets $B_w=\{k : (n{+}1)\wn{wk}<N\}$ cover $\Z_N$. Five empirically
  distinct obstruction classes from the program's earlier experiments
  collapse into \emph{additive orders of one covering phenomenon}
  (Table~\ref{tab:taxonomy}). The framework yields proved lemmas: an exact
  size law (Lemma~\ref{lem:size}), a polygonal witness lemma
  (Lemma~\ref{lem:jgon}), moat counting (Lemma~\ref{lem:counting}), a
  coincidence-counting lemma (Lemma~\ref{lem:D3}), a lift law
  (Lemma~\ref{lem:lift}), and a class-group reduction at prime moduli
  (Lemma~\ref{lem:class}).
\item \textbf{The rigid zone and its classifications}
  (\S\ref{sec:rigid}). At each rung $n$ there is a small set of
  \emph{unit-capable} moduli---those $N$ for which $n-1$ unit bad sets can
  cover $\Z_N$: $\{7,11,13\}$ at $n=4$ and $\{7,13,17,19,37\}$ at $n=5$
  within the verified range. We classify the covering configurations
  completely at every one of these primes
  (Theorems~\ref{thm:K4}--\ref{thm:K37}): each classification is a short
  explicit list of translate classes in a cyclic group (antipodal classes,
  domino tilings, perfect matchings, and at $19$ and $37$ the two classes
  $\{0,1,2,7\}$, $\{0,2,4,6\}$ and the single class $\{0,2,12,15\}$).
  This is the paper's central structural finding: \emph{at rigid primes,
  the covering obstruction is governed by a small translate-class
  structure}. The measured coverage rates of \S\ref{sec:measure} are a
  consequence of this structure, not the primary result.
\item \textbf{The measurement} (\S\ref{sec:measure}). On all
  $3{,}712{,}576$ primitive $5$-sets with speeds $\le 56$: proved lemmas
  close $75.6\%$ ($78.2\%$ with the coincidence lemma), per-modulus
  verification closes $100\%$, and every set has a certifying pair at a
  never-covering configuration. The proved-coverage rate declines as the
  corpus grows ($93.4\%\to78.2\%$ over $V=24\to56$ on the augmented
  battery, crossing below the continuation threshold at $V=56$); the trend
  does not establish stabilization, and the decline is reported as a
  finding, with the residual localized to one named class (mixed order
  signatures at composite moduli).
\item \textbf{The conditional scaling-closure theorem}
  (\S\ref{sec:SC}, \S\ref{sec:scale}). The frontier cells
  $(1K,T{-}3U)$ at $N=p^2$---beyond the reach of enumeration---are
  attacked structurally: the fiber cascade of \S\ref{sec:higherrung}
  reduces, exactly, to iterated dilate-intersections of census solution
  sets (Lemma~\ref{lem:twofiber}), and the closure question becomes a
  pool-versus-cut arithmetic. What is proved: the pool law, the spacing
  lemma, the $k$-cancellation, the free reflection fiber
  (Lemma~\ref{lem:reflect}), the chain-class dilate theorem
  (Theorem~\ref{thm:chaindilate}), and the renormalization lemma
  (Lemma~\ref{lem:renorm}). What is measured: the pool saturation at two
  rungs (the $840$ identity is a $p=17$ accident, disclosed as such),
  the zone-resolved cut quality, and the $T=8$ volume law with margin
  growing in $k$. What remains is consolidated into \emph{one} named
  conditional theorem (Theorem~\ref{thm:SC}) with a six-entry hypothesis
  ledger, a free-parameter disclosure ($c_1$ data-bounded, $\Lambda_0\approx4$
  data-chosen, $\gamma_0$ data-bounded, the fitted $R$ retired), and
  sampled estimates carrying their intervals and \emph{projected} labels.
  The two open obligations are reduced to named objects: the strand
  census (for H1g) and the sharp pair-line constants (for H2m).
  The boundary between those two lists has a structural type, and it
  explains where the paper stops: every proved result in the
  program---everything enumerated above and everything in the sections
  that follow---is a closure, invariance, exhaustiveness, or
  machine-verification statement, while every remaining obligation (the
  H1g margin over the trivial bound, the volume law's growth in $k$, the
  chain class's gap from measured to proved constants) is a uniform
  growth-in-$k$ margin over a trivial or measured baseline. No turn of
  the program has proved a statement of the second type; the audit of
  \S\ref{sec:scale-status} verifies this statement-by-statement,
  including the one candidate that reads like an exception, the cut
  law's ``$k$-cancellation,'' which cancels the $k$-dependence---a
  $k$-invariance, the exact opposite of a growth-in-$k$ margin. The
  empirical half of the reading is thereby settled; the structural half
  (that a classification toolkit cannot produce the missing type) is a
  hypothesis about tool reach. If it is right, the reformulation cannot
  close the conjecture in its current form: it closes the covering
  obstruction exactly where classification suffices, and the residual is
  a growth-in-$k$ margin of a type that has resisted every census-shaped
  attack the program has tried---a structural gap, not a technical one.
  The one margin-shaped candidate was then executed as a proof attempt
  on the committed cells under a pre-committed success criterion
  ($\rho\le 1-\epsilon$, $\epsilon$ independent of $k$), with the state
  space and the whole-cascade comparison target fixed beforehand. It
  closed by theorem for the lossless state space, not by measurement:
  the constrained transfer
  matrices of the lossless cascade are partial permutations
(Lemma~\ref{lem:monomial}), every spectral quantity in the framework
is exactly $0$ or $1$, the cyclic product is diagonal---its trace is
the cascade's own survivor count---and the survival-fraction deficit
the criterion demanded is, in this language, precisely the survivor
inequality Hq and the volume law H1g restated. The instrument is the
census in different notation; it consumes the hypotheses it was
commissioned to prove. That is the mechanism behind the type mismatch:
losslessness forbids generated decay, and the reformulation is bounded
by its own fidelity. (The closure theorem covers automata built on the
cascade's bijective drift---the lossless state space; instruments that
do not inherit the drift are outside its scope, as
\S\ref{sec:scale-status} records.)
\item \textbf{Open problems, and the post-submission structural
  additions} (\S\ref{sec:open}). The pair-selection law, the kernel world,
  the composite micro-cases $49=7^2$, $77=7\cdot 11$, $91=7\cdot 13$,
  rigid-zone termination, and whether the decay is fundamental. The
  post-submission additions carry the composite side to its current
  shape: the group-ring translation and the four-family census
  (\S\ref{sec:groupring}), the $pq$ tensor side (\S\ref{sec:pq}), the
  discharge of the foot-count input (\S\ref{sec:lemmaE}), and---new in
  this revision---the higher rungs (\S\ref{sec:higherrung}): the
  general-$T$ fiber machinery, the $T=8$-critical census with its eleven
  $p$-independent families, the $T=7$ core--tail census, and the shear
  step re-established against growing pinned sets, answering the paper's
  third-rung question affirmatively for the census+shear method; and,
  completing this revision, the scaling-closure record of
  \S\ref{sec:scale}---the derivation, calibration, and epistemic status
  behind Theorem~\ref{thm:SC}. The final revision records the disposition
  of the third bridge of \S\ref{sec:open}: the zonotope--Ehrhart
  instrument engaged as a full strand and closed by the same type
  mismatch (stable arithmetic, false geometric reading), with the
  negative half of the record---the falsified strengthenings, the exact
  witnesses, and the certification ladder---published in the companion
  paper \cite{companion-paper}.
\end{enumerate}

\subsection{What this paper does not claim}

Because the subject is a famous open problem, the boundary of the claims
is stated as carefully as the claims themselves. The paper does not solve,
reduce, or weaken the Lonely Runner Conjecture; the equivalence of
Corollary~\ref{cor:tau} is exact, so the full difficulty of
$\mathrm{LRC}$-$_n$ is present in $(\mathrm{TAU}$-$_n)$. It proves no new
cases of the conjecture at any $n$; the rung results of
\cite{Rosenfeld8,Trakulthongchai,Sungkawichai,Allikvere} are untouched and
unaffected. The rigid-zone map of \S\ref{sec:rigid} is an empirical
finding about a verified range (primes $\le 150$, composites $\le 111$),
not a theorem of finiteness; the successor conjecture
(Conjecture~\ref{conj:successor}) is stated as a conjecture. The coverage
measurements of \S\ref{sec:measure} are measurements; their scaling
beyond the tested range is explicitly open. Finally, the analytic core of
the equality theorem is folklore \cite{Tao-blog}; only the sharpening, the
exact form, and everything built on them are claimed as new, and even
there the correct prior, given how actively the spectrum line computes
exact optima \cite{Cordella}, is that pieces of this may exist in another
language somewhere in the literature we could not reach.

One boundary statement is specific to the conditional theorem.
Theorem~\ref{thm:SC} is exactly what its name says: its two substantive
hypotheses (H1g and the H2 pair) are open, are not established anywhere
in this paper, and no consequence that requires them is claimed
unconditionally. If both close, the reformulation yields a uniform
in-$T$ closure of the $p^2$ covering obstruction and hence, by the
lossless equivalence, the conjecture at the covered rungs; if either
fails, the reformulation's scope is precisely the bounded statement of
Theorem~\ref{thm:SC}---a conditional theorem with named hypotheses, not
a proof. The measured calibration ($T=8,9,10$ at $p\le73$; committed
ground-truth cells at $p\le61$) is stated with its ranges and is not
extrapolated beyond them except where explicitly labeled
\emph{projected}.

\subsection{Notation}

Throughout, $n\ge 2$ is the number of speeds, $T=n+1$ is the loneliness
target denominator, and $V$ has $\gcd(V)=1$ (dividing by the gcd does
not change $M$). For a modulus $N$ we write $\wn{x}=\min(x\bmod N,\;
N-(x\bmod N))$, $\ord_N(w)=N/\gcd(w,N)$ for the additive order of $w$
in $\Z_N$, and $h=\floor{(N-1)/T}$ for the covering threshold index
of \S\ref{sec:cover}. The letter $k$ always denotes a lattice time
index. Computations are exact integer/rational arithmetic wherever a
theorem or table entry asserts an exact value; the Wilson intervals,
the fitted $\Lambda_0$, the projected pools of Table~\ref{tab:pool},
and displayed asymptotics are statistical or asymptotic estimates and
are labeled as such in the free-parameter disclosure. The verification
record is the subject of \S\ref{sec:sound}.
